\section{Discussion}
\label{sec:conclusion}
The upcoming subgoal is accessible before its implementation through a readout of selected attention-head contributions. These contributions can also steer code in another language, and controlled restrictions establish causal involvement of a broader selected set in native implementation. Query and task-span interventions show that the representation interacts with context, while OVReader traces its accessibility through model depth.

\paragraph{Meaning and realization.}
Cross-language implementation transfer supports a representation that can guide more than one surface realization. Natural-language verbalization provides a complementary description of its content. Neither result establishes that the internal representation is itself natural language, or that every detail of the eventual implementation has been selected at the capture point. The position-wise evidence concerns an upcoming operation, not a complete long-horizon plan.

\paragraph{Causal involvement and selection.}
\label{sec:mechanism:competence}
The ablations distinguish effective readout from native dependence on the selected heads. They do not establish a unique bottleneck or demonstrate adaptive rerouting. Similarly, a query that favors a source goal and a KV intervention that changes the executed goal do not identify a unique progress representation. Establishing the algorithm of subgoal selection requires contrasts that independently vary completed work, task order, and target identity while controlling intervention sites.

\paragraph{Scope and limitations.}
Our principal quantitative setting is controlled code generation with a known frontier. Oracle agreement can credit partial implementations and does not certify executable correctness. Readout, task execution, and successful transfer are therefore distinct outcomes. The five capture positions are repeated observations on each specification, and the head-budget selections are not strictly nested. Cross-model results replicate the procedure within each model; they do not show interchangeability of latent coordinates. Free-form code and non-code examples remain exploratory. Boundary-signal interventions are reported as a separate supporting direction in \Cref{app:boundary}.

The results provide a concrete way to examine the semantic operation a model is about to pursue. They motivate further work on when these readouts fail, how subgoals are selected without explicit decomposition, and how their representations change as computation proceeds.
